\section{Compton Streuung}

\subsection*{Ringspektroskopie}
\begin{center}
    \includegraphics*{skizze1.pdf}
\end{center}
Über die Ringspektroskopie können wir die Energie der gestreuten Photonen bei einem bestimmten Winkel $\theta$ bestimmen. Dazu nutzen wir den Abgebildeten Aufbau und verstellen $a$, $b$ und $r$ so, dass sich die gewollten Winkel über
\begin{equation*}
    \theta = \arctan\left(\frac{d/2}{a}\right) + \arctan\left(\frac{d/2}{b}\right)
\end{equation*}
ergibt. Dabei mitteln wir den Strahlengang der Photonen über Streuung in der Mitte des Randes des Rings sowie Detektion in der Mitte des Szintillators.\\
\\
Die Streckenmessung für $a$ und $b$ erfolgt über Skalen auf den Schienen sowie einem Geodreieck.\\
Für $a$: Der Szintillator wird über zwei Stützen gehalten und wir lesen den Ort $a_1$ der Stütze des Kopfes ab. Zusätzlich messen wir den Abstand $a_2$ zwischen Mitte der Stütze und Anfang des Szintillators. Zudem lesen wir den Ort $a_3$ der Mitte der Stütze des Ringes ab. Damit ergibt sich $a = |a_3 - a_1 - a_2|$.\\
Für $b$: Der Gammastrahler ist in Richtung des Ringes auf eine Halterung aufgesteckt. Diese ragt also um die Strecke $b_1$ aus der Halterung raus und wir messen diese Strecke. Die Halterung ist auf einer Stütze, dessen Ort $b_2$ wir ablesen sowie messen wir den Abstand $b_3$ zwischen Mitte der Halterung und Rand. Damit ergibt sich $b = |b_2 - a_3 - b_1 - b_3|$.\\
Für $d$: Wir messen den Durchmesser des gesamten Ringes $d_1$ und die Dicke des Randes $d_2$. Für den Durchmesser eines gedachten Ringes durch die Mitte des ausgedehnten Ringes ziehen wir also 2 mal die Hälfte der Dicke des Randes ab. Damit ergibt sich $d = d_1 - d_2$.\\
Während des Versuches verändern wir nur die Position des Szintillators $a_1$ und die Position der Halterung des Gammastrahlers $b_2$ um die gewünschten Winkel zu erreichen. Alle anderen Werte werden konstant gehalten, sodass für jede Messung die folgenden Werte gelten
\begin{center}
    \begin{tabular}{|c|c|}
        \hline
        \textbf{Variable} & \textbf{Wert}\\
        \hline
        $a_2$ & $\SI{1.0 \pm 1.0}{\centi\meter}$ \\
        \hline
        $a_3$ & $\SI{160.0 \pm 1.0}{\centi\meter}$ \\
        \hline
        $b_1$ & $\SI{3.5 \pm 1.0}{\centi\meter}$ \\
        \hline
        $b_3$ & $\SI{3.9 \pm 1.0}{\centi\meter}$ \\
        \hline
        $d_1$ & $\SI{25.0 \pm 1.0}{\centi\meter}$ \\
        \hline
        $d_2$ & $\SI{1.5 \pm 1.0}{\centi\meter}$ \\
        \hline
        $d$ & $\SI{23.5 \pm 1.0}{\centi\meter}$ \\
        \hline
    \end{tabular}
\end{center}

Für die Unterschiedlichen Winkel $\theta$ Stellen wir die Folgenen Werte für $a_1$ und $b_2$ ein
\begin{center}
    \begin{tabular}{|c|c|c|}
        \hline
        $\theta_{\text{Ziel}}$ & $a_1$ & $b_2$ \\
        \hline
        $\SI{10}{\degree}$ & $\SI{19.0 \pm 1.0}{\centi\meter}$ & $\SI{296.7 \pm 1.0}{\centi\meter}$ \\
        \hline
        $\SI{20}{\degree}$ & $\SI{102.0 \pm 1.0}{\centi\meter}$ & $\SI{246.7 \pm 1.0}{\centi\meter}$ \\
        \hline
        $\SI{30}{\degree}$ & $\SI{124.5 \pm 1.0}{\centi\meter}$ & $\SI{226.7 \pm 1.0}{\centi\meter}$ \\
        \hline
        $\SI{40}{\degree}$ & $\SI{131.8 \pm 1.0}{\centi\meter}$ & $\SI{206.7 \pm 1.0}{\centi\meter}$ \\
        \hline
        $\SI{50}{\degree}$ & $\SI{137.0 \pm 1.0}{\centi\meter}$ & $\SI{196.7 \pm 1.0}{\centi\meter}$ \\
        \hline
    \end{tabular}
\end{center}
Es ergeben sich damit die folgenden Werte für $a$, $b$ und $\theta$
\begin{center}
    \begin{tabular}{|c|c|c|}
\hline
$a$ & $b$ & $\theta_{\text{Messung}}$ \\
\hline
$\SI{140.0 \pm 1.7}{\centi\meter}$ & $\SI{129.3 \pm 2.0}{\centi\meter}$ & $\SI{10.0 \pm 0.6}{\degree}$ \\
\hline
$\SI{57.0 \pm 1.7}{\centi\meter}$ & $\SI{79.3 \pm 2.0}{\centi\meter}$ & $\SI{20.1 \pm 1.2}{\degree}$ \\
\hline
$\SI{34.5 \pm 1.7}{\centi\meter}$ & $\SI{59.3 \pm 2.0}{\centi\meter}$ & $\SI{30.0 \pm 1.9}{\degree}$ \\
\hline
$\SI{27.2 \pm 1.7}{\centi\meter}$ & $\SI{39.3 \pm 2.0}{\centi\meter}$ & $\SI{40.0 \pm 2.6}{\degree}$ \\
\hline
$\SI{22.0 \pm 1.7}{\centi\meter}$ & $\SI{29.3 \pm 2.0}{\centi\meter}$ & $\SI{50.0 \pm 3.3}{\degree}$ \\
\hline
\end{tabular}
\end{center}